\subsection{实际执行行为的补充分析}\label{sec:supporting}\label{sec:f1-collapse}
MATH 对照表明，结果不同未必需要程序不同。补充的 BIRD 研究考察反方向的问题：不同源码是否真正产生不同且有用的执行行为？这里的机制指已实现的提示变换、工具使用或控制流操作。该研究使用独立总体、评分及执行协议，不能解释 MATH 结果的成因。

\paragraph{代码差异与观测结果。}
所有 $12$ 个发现候选都通过代码差异检查，但 $66$ 个生成候选对中有 $21$ 对的正确性向量相同；计入基线时为 $36/91$，源码相似度与正确性分歧的秩相关为 $\rho=-0.010$。这些 $60$ 题上的一致不等于一般程序等价，但说明需要直接检查执行（附录图~\ref{fig:phenomenon}）。

\paragraph{轨迹定位三种不同的失败。}
\textbf{T1：描述了但没有实现。}$12$ 个候选中的 $6$ 个在文档中宣称检索或自检，实际只有一次模型调用且不执行产物。\textbf{T2：存在但未触发。}全部 $10$ 道已审 \react{} 任务的第一次 SQL 均执行成功，修复分支从未触发，最终 SQL 与 \bare{} 相同；执行成功不代表语义正确。\textbf{T3：执行了但输出失败。}一个候选执行多轮修复，却把散文作为最终 SQL，丢失了基线答对的任务。三者分别要求实现检查、能触发分支的任务和有效的输出合同；它们不确立激活机制可以提高准确率。附录~\ref{app:bird-traces} 保留完整样例、人工控制及其开发暴露限制。

独立的 BIRD 准入研究还发现，强制策略加门控相对自由无门控的平均余量差为 $\RevisionPrimaryDelta$\,\%。该门控捆绑了结构排除与合规验证，不能识别纯验证效应（附录~\ref{app:phase2-protocol}）。因此，源码与准入证据需要通过实际执行来核实，其任务优势及选择价值则仍需分别测量。
